\documentclass[profile=standard]{ipara-handbook}
\title{B10｜样本中心化与正交投影：自由度与抽样分布}
\author{}
\date{}
\begin{document}
\maketitle

\begin{quote}
本册回答：\textbf{为什么样本方差偏偏除以 $n-1$？为什么正态总体下 $\bar X$ 与 $S^2$ 会独立？为什么离差平方和会出现 $\chi^2(n-1)$？这三件事其实来自同一张正交分解图。}
\end{quote}

\section{本册定位：把统计样本看成欧氏空间中的随机向量}
概率 Topic07 Own 样本均值、样本方差、三大抽样分布；B00/线代 Topic01 Own 正交投影；线代 Topic03 Own 子空间维数与秩。B10 只 Own 下面这条接口：
\[
\boxed{
\mathbb R^n
=\operatorname{span}\{\mathbf1\}
\oplus
\mathbf1^\perp
}
\]
以及样本向量在这两个正交子空间中的分解如何翻译成“均值信息 + 残差信息”。

\section{先固定样本对象：均值、残差与样本方差}
令样本向量
\[
X=(X_1,\dots,X_n)^T,
\qquad
\mathbf1=(1,\dots,1)^T,
\]
并定义
\[
\bar X:=\frac1n\sum_{i=1}^nX_i,
\qquad
S^2:=\frac1{n-1}\sum_{i=1}^n(X_i-\bar X)^2.
\]
前半册的投影恒等式纯属线性代数，对 $X_i$ 的分布没有要求；只有在后面讨论“无偏”“独立”“$\chi^2$ 与 $t$”时，才逐步加入独立同分布、有限方差或正态性等概率条件。这样可以避免把几何恒等式与概率定理混成同一个前提包。

\section{样本均值与均值方向的投影是什么关系}
投影到 $\operatorname{span}\{\mathbf1\}$ 的矩阵为
\[
P=\frac{\mathbf1\mathbf1^T}{\mathbf1^T\mathbf1}
=\frac1n\mathbf1\mathbf1^T.
\]
于是
\[
PX
=\frac1n\mathbf1\mathbf1^TX
=\bar X\mathbf1.
\]
所以 $\bar X\mathbf1$ 是样本向量在“所有坐标共同移动”方向上的投影向量。若改用单位向量
\[
e_1:=\frac1{\sqrt n}\mathbf1,
\]
则真正的规范正交坐标为
\[
e_1^TX=\sqrt n\,\bar X.
\]
因此必须区分：$\bar X$ 是样本平均值，$\sqrt n\,\bar X$ 是单位均值方向上的坐标，$\bar X\mathbf1$ 才是投影向量。三者相关，但不是同一个对象。

\section{中心化 = 删除均值方向}
定义残差投影矩阵
\[
R:=I-P.
\]
则
\[
RX=X-\bar X\mathbf1
=(X_1-\bar X,\dots,X_n-\bar X)^T.
\]
先定义均值方向的正交补
\[
\mathbf1^\perp:=\{r\in\mathbb R^n:\mathbf1^Tr=0\}.
\]
由于
\[
\mathbf1^T RX=0,
\]
中心化后的残差向量落在
\[
\boxed{\mathbf1^\perp}
\]
中。

这就是恒等式
\[
\sum_{i=1}^n(X_i-\bar X)=0
\]
的线代意义：$n$ 个残差坐标不是 $n$ 个可自由变化的坐标，它们被一个线性约束限制在 $n-1$ 维子空间中。这里先谈的是几何维数，不把“线性约束”偷换成概率独立性结论。

\section{为什么自由度是 $n-1$}
投影矩阵 $P$ 的秩为 $1$，因此
\[
\operatorname{rank}(R)
=\operatorname{rank}(I-P)
=n-1.
\]
而残差平方和就是残差向量的平方范数：
\[
\sum_{i=1}^n(X_i-\bar X)^2
=\lVert RX\rVert^2
=X^TRX,
\]
因为 $R=R^T=R^2$。

所以“$n-1$ 个自由度”的核心不是一句统计口诀，而是：
\[
\boxed{\text{中心化投影删掉了一个均值方向，残差只剩 }n-1\text{ 个正交方向。}}
\]

按上面的定义，样本方差写成
\[
S^2=\frac{1}{n-1}\lVert RX\rVert^2.
\]
几何上，分母 $n-1$ 与残差空间维数对应；概率上，若 $X_1,\dots,X_n$ 独立同分布且方差为 $\sigma^2<\infty$，还进一步有
\[
E\lVert RX\rVert^2=(n-1)\sigma^2,
\qquad
E[S^2]=\sigma^2.
\]
这条期望关系也能由 B09 的协方差矩阵接口一行生成。令 $Y:=X-\mu\mathbf1$，则 $\operatorname{Cov}(Y)=\sigma^2I$，而 $RY=RX$。记 $\operatorname{tr}(A)$ 为矩阵对角元之和，则
\[
\begin{aligned}
E\lVert RX\rVert^2
&=E(Y^TRY)\\
&=\operatorname{tr}\!\left(R\,E[YY^T]\right)\\
&=\sigma^2\operatorname{tr}(R)\\
&=(n-1)\sigma^2,
\end{aligned}
\]
因为正交投影 $R$ 的特征值只有 $0$ 与 $1$，且 $\operatorname{rank}(R)=\operatorname{tr}(R)=n-1$。所以“维数 $n-1$”解释自由度来源，而“除以 $n-1$ 得到无偏估计”还需要随机样本模型，不能只靠几何一句话包办。

\section{正态样本为什么把几何分解升级成概率独立}
设样本相互独立且同分布，本册缩写为 i.i.d.，并且
\[
X_i\overset{\mathrm{i.i.d.}}\sim N(\mu,\sigma^2).
\]
标准化样本向量
\[
Z:=\frac{X-\mu\mathbf1}{\sigma}
\]
满足
\[
Z\sim N(0,I_n).
\]
这是各向同性标准正态：任意正交变换后分布不变。

把空间选成一组规范正交基，其中第一根为
\[
e_1=\frac1{\sqrt n}\mathbf1,
\]
其余 $e_2,\dots,e_n$ 张成 $\mathbf1^\perp$。在这组正交坐标中，
\[
e_1^TZ
\]
只描述均值方向，而
\[
e_2^TZ,\dots,e_n^TZ
\]
只描述残差方向。

由于联合标准正态在正交变换后仍是各坐标独立的标准正态，这两块正交坐标彼此独立。

因此在正态样本下：
\[
\boxed{PX\text{ 与 }RX\text{ 相互独立}}
\]
进一步得到样本均值与样本方差概率独立。本册沿用概率总图的记号约定：双竖 $\perp\!\!\!\perp$ 表示概率独立，单竖 $\perp$ 只表示几何正交，因此
\[
\boxed{\bar X\perp\!\!\!\perp S^2.}
\]

这里必须强调：一般随机向量的几何正交\textbf{不推出}概率独立；正态结构是关键加强条件。

\section{卡方自由度为什么就是残差空间维数}
先固定分布定义：若 $Z_1,\dots,Z_\nu$ 相互独立且都服从 $N(0,1)$，则
\[
U:=\sum_{j=1}^{\nu}Z_j^2
\]
称为服从自由度 $\nu$ 的卡方分布，记作 $U\sim\chi^2(\nu)$。

在上面的规范正交坐标中，残差向量对应 $n-1$ 个独立标准正态坐标：
\[
Z_2,\dots,Z_n\overset{\mathrm{i.i.d.}}\sim N(0,1).
\]
所以
\[
\frac{\lVert RX\rVert^2}{\sigma^2}
=Z_2^2+\cdots+Z_n^2
\sim\chi^2(n-1).
\]
于是
\[
\boxed{
\frac{(n-1)S^2}{\sigma^2}
\sim\chi^2(n-1).
}
\]

因此在这个正态样本模型里，“自由度 $n-1$”在三个层面严格对齐：
\begin{itemize}
\item 残差子空间的维数；
\item 投影矩阵 $R$ 的秩/迹；
\item 卡方中独立标准正态平方项的个数。
\end{itemize}

\section{$t$ 分布为什么自然从两块正交信息拼出来}
先固定定义：若 $Z\sim N(0,1)$、$U\sim\chi^2(\nu)$ 且二者概率独立，则
\[
T:=\frac{Z}{\sqrt{U/\nu}}
\]
服从自由度 $\nu$ 的 $t$ 分布。

正态样本下
\[
\frac{\sqrt n(\bar X-\mu)}{\sigma}\sim N(0,1),
\]
而
\[
\frac{(n-1)S^2}{\sigma^2}\sim\chi^2(n-1),
\]
且二者独立。

因此
\[
\frac{\sqrt n(\bar X-\mu)}{S}
=\frac{Z}{\sqrt{U/(n-1)}}
\sim t(n-1).
\]

所以 $t$ 分布不是“把未知 $\sigma$ 换成 $S$ 后神奇地产生的新表”，而是\textbf{均值方向的标准正态坐标 ÷ 独立残差方向产生的随机尺度}。

这是一条非常重要的结构生成链：
\[
\boxed{\text{正交分解}\to\text{独立}\to\chi^2\text{ 残差尺度}\to t\text{ 标准化}.}
\]

\section{平方和分解也是勾股定理}
对任意固定常数 $\mu$，
\[
X-\mu\mathbf1
=(X-\bar X\mathbf1)+(\bar X-\mu)\mathbf1.
\]
两项正交，因此由勾股关系
\[
\boxed{
\sum_{i=1}^n(X_i-\mu)^2
=\sum_{i=1}^n(X_i-\bar X)^2
+n(\bar X-\mu)^2.
}
\]

这条经典统计平方和分解，本质上就是 B00 的正交投影分解在样本空间中的实例。

\section{母例一：三个观测怎样把“中心化”画成一个二维残差平面}
先不谈概率，只看一次已经观察到的数据
\[
x=(1,2,6)^T.
\]
样本均值为
\[
\bar x=3,
\]
所以均值投影与残差分别是
\[
Px=(3,3,3)^T,
\qquad
Rx=(-2,-1,3)^T.
\]
检查
\[
\mathbf1^TRx=-2-1+3=0,
\]
说明残差确实落在 $\mathbf1^\perp$ 中。虽然残差写了三个坐标，但第三个已经被前两个通过“总和为零”决定，所以真正只剩二维自由度。

残差平方和为
\[
\lVert Rx\rVert^2=4+1+9=14,
\]
因此按本册样本方差定义
\[
s^2=\frac{14}{3-1}=7.
\]
再用单位均值方向 $e_1=\mathbf1/\sqrt3$，其投影坐标为
\[
e_1^Tx=3\sqrt3,
\]
再次提醒：$3$ 是样本均值，$3\sqrt3$ 才是单位方向坐标，而 $(3,3,3)^T$ 是投影向量。

\section{母例二：$n=2$ 时一眼看见“正交 $\to$ 独立 $\to\chi^2$”}
设
\[
X_1,X_2\overset{\mathrm{i.i.d.}}\sim N(\mu,\sigma^2).
\]
对标准化样本做正交坐标变换：
\[
Z_+=\frac{X_1+X_2-2\mu}{\sigma\sqrt2}
=\frac{\sqrt2(\bar X-\mu)}{\sigma},
\]
\[
Z_- =\frac{X_1-X_2}{\sigma\sqrt2}.
\]
这两个方向分别是 $\mathbf1$ 方向和其正交补方向。因为原标准化样本是二维联合标准正态，而变换矩阵正交，所以
\[
Z_+,Z_-\overset{\mathrm{i.i.d.}}\sim N(0,1).
\]
特别地，它们概率独立。

另一方面，$n=2$ 时
\[
\sum_{i=1}^2(X_i-\bar X)^2
=\frac{(X_1-X_2)^2}{2},
\]
故
\[
\frac{S^2}{\sigma^2}
=Z_-^2
\sim\chi^2(1).
\]
而 $\bar X$ 只由 $Z_+$ 决定，$S^2$ 只由 $Z_-^2$ 决定，因此
\[
\bar X\perp\!\!\!\perp S^2.
\]

这个 $n=2$ 母例把抽象结论压到最小维数：一个均值方向 + 一个残差方向。自由度 $1$、卡方平方项个数 $1$、均值与残差独立三件事完全重合。

\section{B10 串起哪些章节}
\begin{tabular}{p{0.23\linewidth}p{0.30\linewidth}p{0.39\linewidth}}
\hline
Owner / 下游 & 提供对象 & B10 的翻译责任\\
\hline
B00 / 线代 Topic01 & 正交投影、勾股分解 & 把样本均值/中心化翻译成两块正交分量\\
线代 Topic03 & rank、子空间维数 & 把自由度读成残差投影的秩 $n-1$\\
概率 Topic07 & $\bar X,S^2,\chi^2,t,F$ & 解释卡方自由度与均值/方差独立的结构来源\\
概率 Topic04/B09 & 正态向量与协方差 & 提供正交变换下 Gaussian 独立性的加强条件\\
\hline
\end{tabular}

\section{为什么“估计一个参数就减一”不是根本规则}
在本问题中减一，是因为我们把样本投影到一个一维均值子空间，并把这一维从残差空间中删掉。

更一般的线性模型中，若投影到一个 $p$ 维模型子空间，残差空间维数会变成 $n-p$。这说明“自由度减少多少”真正由\textbf{消耗了多少独立线性方向}决定。

一般 Cochran 定理、回归与方差分析属于 Extension；数学一只保留当前一维均值投影这条主线。

\section{反桥梁}
\begin{tabular}{p{0.20\linewidth}p{0.04\linewidth}p{0.20\linewidth}p{0.48\linewidth}}
\hline
A & & B & 真正区别\\
\hline
几何正交 & $\neq$ & 概率独立 & 只有联合正态等额外结构下才能从正交升级为独立\\
残差有 $n$ 个坐标 & $\neq$ & 有 $n$ 个自由度 & 它们满足总和为零，只张成 $n-1$ 维子空间\\
除以 $n-1$ & $\neq$ & 任意无偏修正口诀 & 本质来自当前残差空间维数；其他模型自由度可能不同\\
$\chi^2(n-1)$ & $\neq$ & 任意总体下精确成立 & 精确卡方抽样分布依赖正态总体\\
\hline
\end{tabular}

\section{调用信号与一页压缩}
看到“为什么样本方差除 $n-1$、$\bar X$ 与 $S^2$ 独立、正态样本卡方、$t$ 分布自由度”时，先画一张空间分解：
\[
\boxed{
X
=PX+RX
=\bar X\mathbf1+(X-\bar X\mathbf1),
\qquad
\operatorname{span}\{\mathbf1\}\perp\mathbf1^\perp.
}
\]

然后问：
\begin{enumerate}
\item 均值方向占几维？
\item 残差空间还剩几维？
\item 当前总体是否正态，足以把正交分量升级为独立？
\item 当前平方和是否正好是残差投影的平方范数？
\end{enumerate}

能从这张图重新生成 $n-1$、$\chi^2$、独立性和 $t$，就不再需要把它们当四条孤立定理背诵。

\end{document}
